\let\negmedspace\undefined
\let\negthickspace\undefined
\documentclass[journal,12pt,onecolumn]{IEEEtran}
\usepackage{cite}
\usepackage{amsmath,amssymb,amsfonts,amsthm}
\usepackage{algorithmic}
\usepackage{graphicx}
\graphicspath{{./figs/}}
\usepackage{textcomp}
\usepackage{xcolor}
\usepackage{txfonts}
\usepackage{listings}
\usepackage{enumitem}
\usepackage{mathtools}
\usepackage{gensymb}
\usepackage{comment}
\usepackage{caption}
\usepackage[breaklinks=true]{hyperref}
\usepackage{tkz-euclide} 
\usepackage{listings}
\usepackage{gvv}                                        
\usepackage{amsmath}                                 
\usepackage[latin1]{inputenc}     
\usepackage{xparse}
\usepackage{color}                                            
\usepackage{array}                                            
\usepackage{longtable}                                       
\usepackage{calc}                                             
\usepackage{multirow}
\usepackage{multicol}
\usepackage{hhline}                                           
\usepackage{ifthen}                                           
\usepackage{lscape}
\usepackage{tabularx}
\usepackage{array}
\usepackage{float}

\begin{document}

\title{2.7.10}
\author{AI25BTECH11038 -- Tejas Uppala}
{\let\newpage\relax\maketitle}

\textbf{Question:}

\noindent Find the direction and normal vectors of the following line $2x + 3y = 9$

\textbf{Solution:}

\noindent For a general line y = mx + c, it can be converted into its vector form in the following way,

\begin{equation}
    y = mx + c
\end{equation}

\begin{equation}
    \myvec{x\\y} = \myvec{x\\mx+c} = \myvec{0\\c} + x\myvec{1\\m}    
\end{equation}

\noindent yielding

\begin{equation}
    \textbf{x} = \textbf{h} + k\textbf{m}
\end{equation}

\noindent where \textbf{h} is any point on the line and \textbf{m} is the direction vector of the line,

\begin{equation}
    \textbf{m} = \myvec{1\\m}
\end{equation}

\noindent The given line can be written as, 

\begin{equation}
    y = \brak{\frac{-2}{3}}x + \frac{9}{3} = \brak{\frac{-2}{3}}x + 3
\end{equation}

\noindent on correlating,

\begin{equation}
    \textbf{m} = \myvec{1 \\ \frac{-2}{3}} or \myvec{3\\-2}, \textbf{h} = \myvec{0\\3}
\end{equation}

\noindent The \textbf{normal vector} (n) of the line is \myvec{-m\\1}. Therefore, 

\begin{equation}
    \textbf{n} = \myvec{\frac{2}{3}\\1} or \myvec{2\\3}
\end{equation}

\end{document}
